\section{讲座定位与问题设置}

这节课的核心，不是给出一个新的 Agent 框架名字，而是回答一个更基础的问题：在 LLM 已经具备 Tool Use、Chain-of-Thought、长上下文和代码执行能力之后，为什么还必须系统性地引入 Multi-Agent 训练视角？Oriol Vinyals 的答案非常明确：如果只做单 Agent 指令优化，系统几乎必然在真实开放环境中出现脆弱性；只有把``协作者''和``对抗者''都纳入训练生态，模型才能获得可持续鲁棒性。

\begin{importantbox}{本讲的主命题}
课程主线可压缩为一句话：\textbf{把 AlphaStar 时代对 population training、league dynamics、exploiter discovery 的经验，迁移到 LLM Agent 的 post-training 时代。}
\end{importantbox}

\begin{knowledgebox}{讲者给出的历史坐标}
本讲并非从零开始设计 LLM Agent，而是沿着 Atari $\rightarrow$ Go $\rightarrow$ StarCraft II $\rightarrow$ LLM Agents 的连续谱系展开。核心是：问题复杂度每上一个台阶，训练目标都会从``单一最优策略''转向``多策略生态中的稳态''。
\end{knowledgebox}

讲者在 18 分钟附近的表述非常关键：\textit{``This is strikingly close to tool use in general.''}\footnote{来源：\href{https://www.youtube.com/watch?v=ntjOxjZMaac}{Agentic AI MOOC Lecture 03}，字幕区间 00:17:49--00:18:05。} 这里的 ``this'' 指的是 AlphaStar 中把高维决策拆成 API-like 动作调用序列的机制。换言之，今天 LLM 在 MCP / function calling 里做的事情，在游戏智能里早已出现过结构同构。

\begin{warningbox}{不要把本讲理解成``游戏经验照搬''}
讲者也反复强调差异：StarCraft 的奖励更清晰、动作接口更结构化，而真实世界 Prompt 与用户目标更含糊。迁移是``借鉴动力学和训练组织方式''，不是``复用原任务定义''。
\end{warningbox}

\subsection{本章小结}
本章给出全讲的评价标准：后续所有技术点都围绕同一目标，即如何让 LLM Agent 从``单点正确''转向``群体互动下的长期稳定正确''。这一目标决定了我们必须讨论 API 设计、IL+RL 混合训练、对抗发现机制和 league 匹配策略。

